% R32 section 3: methods (input by R32_floquet_spectra.tex)
\section{Methods}
\label{sec:methods}

\subsection{Model and conventions}
The sixteen unitary rules are the tuples $(r_{00},r_{01},r_{10},r_{11})\in\{I,V\}^4$, Wolfram numbers
51, 54, 57, 60, 99, 102, 105, 108, 147, 150, 153, 156, 195, 198, 201 and 204,
with $V=H$ throughout this report (Sec.~\ref{sec:x} for $V=X$). Boundary
conditions are open with fixed vacuum neighbours (obc0). One period is
$U=L_BL_A$, with $L_A$ the gates on the odd (1-based) sites, applied first, and
$L_B$ those on the even sites. The quasienergy convention is
$U\ket{\phi}=e^{-i\theta}\ket{\phi}$, $\theta\in(-\pi,\pi]$. The bipartition is
$A=$ sites $1\ldots\lfloor N/2\rfloor$.

\subsection{Pipeline and independent cross-checks}
The diagonalisation is done in Julia by \fn{HSF/generate\_spectral\_data.jl}, an
optimised successor of the earlier \fn{generate\_EE\_data.jl}, which is left
untouched. The analysis is done in Python, in \fn{code/qca\_fragmentation/spectral/}.
\begin{itemize}
\item \textbf{Two independent constructions of $U$} agree. The Julia operator
  is built from Kronecker products of local gates. The Python operator comes
  from this repository's exact $\mathbb{Z}[1/\sqrt2]$ engine
  (\fn{core/cycle.py}). Their quasienergy multisets agree to $<10^{-14}$ for all
  16 rules at $N=8$ (test \fn{test\_engine\_matches\_julia\_spectrum}, $N=6,7,8$).
  Comparisons of angle multisets cut the circle in the middle of the largest
  gap. Sorting raw angles is unsafe because a level at $\theta=0$ computed as
  $-10^{-16}$ wraps to $2\pi$. A first version of the test fell into exactly
  this trap and reported a spurious disagreement.
\item \textbf{Against the stored HSF eigendata} (complex \fn{eigen} per sector):
  quasienergies agree to $10^{-14}$. The entropies agree to $10^{-15}$ for every
  globally non-degenerate level, the only case where the old entropies are
  basis independent.
\end{itemize}

\subsection{Diagonalisation (details and optimisation in Sec.~\ref{sec:algorithm})}
$U$ is real orthogonal (Sec.~\ref{sec:lit}). Each block is diagonalised by
\begin{enumerate}
\item solving the real symmetric problem $C=(A+A^T)/2$, whose eigenvalues are
  $\cos\theta$, with LAPACK's MRRR solver ($O(n)$ workspace);
\item grouping the $\cos\theta$ into clusters (tolerance $10^{-9}$) and
  diagonalising the antisymmetric part $(A-A^T)/2$ inside each cluster, which
  separates $\theta$ from $-\theta$.
\end{enumerate}
This yields exactly orthonormal eigenvectors, also inside degenerate eigenspaces,
which a general complex eigensolver does not guarantee. It is also several times
faster. The three dense products the cluster loop needs are done as single
matrix--matrix multiplications rather than per cluster (that step was memory-bound
and took half the run time). On the 12th-gen i7 used here, dense LAPACK is
fastest with 4--6 BLAS threads and degrades sharply at 20. Every run stores the
largest eigen-residual $\max|U v-e^{-i\theta}v|$ and the orthonormality error.

\subsection{Symmetry resolution}
\label{sec:symres}
Each eigenvector carries three labels, all found and verified at run time:
\begin{enumerate}
\item \textbf{Krylov sector}: connected component of the support of $U$
  (classical fragmentation).
\item \textbf{Gate-algebra block}: the decomposition of each sector by the
  commutant of the bond algebra $\mathcal{A}_0=\langle g_1,\ldots,g_N\rangle$
  generated by the local gates~\cite{moudgalya2022}. This includes quantum
  fragmentation and hidden internal symmetries. It is found without prior
  knowledge of the symmetry:
  \begin{enumerate}
  \item diagonalise a generic element $K=\sum_ia_ig_i$ (random $a_i\in[1,2]$,
    fixed seed);
  \item transform a second generic element $K'=\sum_ib_ig_i$ into that
    eigenbasis;
  \item take the connected components of the transformed $K'$ as the blocks.
  \end{enumerate}
  If the commutant is Abelian, each eigenvector of $K$ lies in one block, and
  $K'$ connects everything within a block. Isomorphic blocks (non-Abelian
  multiplicity) are merged, which is harmless: they give exactly degenerate
  copies of one spectrum, removed in the statistics. \emph{Check}: the largest
  matrix element of $Q^TUQ$ between different blocks is stored
  (\fn{max\_offblock}).
\item \textbf{Spatial label.}
  \begin{itemize}
  \item obc0, odd $N$, mirror-symmetric rules ($r_{01}=r_{10}$): the reflection $R$ commutes with $U$, and every
    eigenvector gets a parity $\pm1$, a true quantum number.
  \item obc0, even $N$: $R$ exchanges the two layers ($RL_AR=L_B$), so $R$ alone is no symmetry. The half step
    $\Omega=RL_A$ is one: it commutes with $U$ and $\Omega^2=U$ (both asserted at run time).
  \item pbc (even $N$ only), every rule: translation by \emph{one} site maps layer $A$ onto layer $B$, so the
    half-translation $\Omega_T=T_1L_A$ commutes with $U$ and $\Omega_T^2=UT_2$, with $T_2$ the translation by two sites
    (asserted). It resolves the $T_2$ momentum $j=0,\ldots,N/2-1$ and splits each momentum sector once more. The
    site-centred reflection does not commute with $\Omega_T$ in general and is not resolved on the ring: a residual
    $\mathbb{Z}_2$ can remain in the $j=0$ and $j=N/4$ sectors of mirror-symmetric rules.
  \end{itemize}
  For the square-root symmetries ($\Omega$, $\Omega_T$) the stored quantity is their eigenphase $\varphi$, obtained by
  diagonalising the symmetry inside each degenerate eigenspace of $U$ (with the same orthonormal $\cos$/$\sin$
  two-step as for $U$). The resolved spectrum of a block is the spectrum of $\varphi$, not of $\theta$: a $\pm1$
  ``sheet'' label derived from $e^{i\theta/2}$ would flip at the branch cut $\theta=\pm\pi$ and mix the two halves
  there. Labels are used for a block only if the block is checked to be invariant under the symmetry, and the invariance
  residual is stored.
\end{enumerate}
Time reversal ($T=L_AK$, $T^2=1$) fixes the reference ensemble to $\beta=1$
(Sec.~\ref{sec:lit}).

\subsection{Statistics}
Gap ratios are computed per fully resolved block (sector, gate-algebra block, true symmetry label), on the block's level
variable ($\varphi$ where a square-root symmetry was resolved, otherwise $\theta$), with exact degeneracies removed.
\begin{itemize}
\item \emph{Folding is decided numerically per block.} If the block spectrum is symmetric under $x\to-x$ (complex
  conjugation maps the block to itself: every block of a real symmetry), only $|x|\in[0,\pi]$ is used. Otherwise, e.g.\
  in a ring-momentum sector $j\ne-j$ whose conjugate partner is another sector, the full circle is used, including the
  wrap-around spacing.
\item \emph{The reference ensemble depends on the block.} Self-conjugate blocks carry the time reversal $T=L_AK$ and are
  compared with $\beta=1$ ($0.5307$). In momentum sectors $j\ne-j$, $T$ maps $j\to-j$, so no antiunitary symmetry is
  left within the sector, and a chaotic sector is $\beta=2$ (CUE, $0.5996$). Both branches are tested against random
  matrices: CUE through the full-circle branch gives $0.597\pm0.005$, and products of random reflections through the
  folding branch give $0.534\pm0.007$.
\item Only blocks with at least 18 ratios are used, and the ratios of all such blocks are pooled. For comparison the
  same statistic is computed at two coarser resolutions: per Krylov sector only, and for the whole spectrum as one list.
\end{itemize}

\subsection{Entanglement and the fair reference}
For every eigenvector the half-chain von Neumann entropy $S_\mathrm{vN}$ is
stored. For every degenerate eigenspace of dimension $g$ (inside its block and
label) the basis-independent Haar-averaged R\'enyi-2 entropy is stored as well,
\begin{equation}
  S_2^\mathrm{Haar}=-\ln\frac{\operatorname{Tr}P_A^2+\operatorname{Tr}P_B^2}{g(g+1)},
  \qquad P_A=\operatorname{Tr}_BP,\ P_B=\operatorname{Tr}_AP,
  \label{eq:s2haar}
\end{equation}
where $P$ is the projector onto the eigenspace. This follows from
$\mathbb{E}[\ket\psi\!\bra\psi^{\otimes2}]=P^{\otimes2}(1+\mathrm{SWAP})/(g(g+1))$,
and for $g=1$ it is the ordinary R\'enyi-2 entropy of the eigenvector. The same
formula with $P$ the projector onto the whole \emph{block} gives
$S_2^\mathrm{rand}(\mathrm{block})$, the entanglement of a random state inside
that block. That is the fair reference for a fragmented system
(Sec.~\ref{sec:lit}), and the ratio
\begin{equation}
  \eta=S_2^\mathrm{Haar}/S_2^\mathrm{rand}(\mathrm{block})
\end{equation}
is $\approx1$ for a thermal eigenstate and $\ll1$ for an atypically weakly
entangled, scar-like one. $\eta$ is defined for blocks of dimension $\ge2$.
States with $\eta<0.5$ in blocks of dimension $\ge20$ are counted as
scar-like.
